\documentclass[conference]{IEEEtran}
\IEEEoverridecommandlockouts

\usepackage{cite}
\usepackage{amsmath,amssymb}
\usepackage{graphicx}
\usepackage{booktabs}
\usepackage{url}
\usepackage[hidelinks]{hyperref}

% ---------------------------------------------------------------------------
% AUTHOR NOTE (delete before submission)
%
% Target venue: LLMSC 2027, International Workshop on LLM Supply Chain Analysis,
% co-located with ICSE 2027. Archival, IEEE format, 4 to 8 pages.
% Deadline 2026-11-27.
%
% DOUBLE BLIND: before submitting, replace the two URLs in Data Availability
% with an anonymous.4open.science mirror, and keep the author block anonymous.
% Restore the real URLs at camera-ready.
%
% Section order: abstract, introduction, related work, taxonomy, tool,
% study design, results, discussion, conclusion.
%
% All numbers come from scripts/08_final_analysis.py, run against
% harness-eval v7.10.1 over the pinned commits in data/corpus_manifest.csv.
% ---------------------------------------------------------------------------

\begin{document}

\title{Auditing the Harness: Composition-Level Defects in\\
AI Code-Agent Configurations}

\author{\IEEEauthorblockN{Anonymous Author(s)}
\IEEEauthorblockA{Affiliation withheld for double-blind review}}

\maketitle

\begin{abstract}
AI code agents such as Claude Code, Cursor, and GitHub Copilot are customized
through persistent configuration artifacts: skills, commands, instruction files,
hooks, MCP servers, and subagent definitions, increasingly installed from public
marketplaces and community repositories. These artifacts, collectively the
\emph{harness}, form a software supply chain surface. They are third-party
instructions that steer agent behavior, execute scripts with developer
privileges, and persist across every session. Unlike conventional dependencies
they lack verification tooling, and their defects produce no error at install
time. Existing analyses inspect these artifacts one file at a time, which leaves
an important class of defect undetectable in principle, because the defect is a
property of how components compose rather than of any component.

We organize harness defects by the scope of analysis their detection requires
and show that this axis is decidable from a detector's implementation, which
makes the boundary of per-file analysis a measurable quantity rather than a
claim. We then instantiate the organization in an open-source auditor
implementing 97 rules across eight assistant platforms, 26 of which require
scope beyond a single file, and apply it to 115 assembled configurations and 65
published skill collections identified from a frame of 5{,}718 repositories. We
find that 85.2\% of real configurations carry at least one defect and 62.6\%
carry an error-severity defect. Most notably, 57.4\% carry a defect that
single-file analysis cannot reach: unreachable skills, circular reference
chains, permission escalation along delegation edges, divergence between the
context files of different assistants, and, in two configurations, a
credential-to-network delegation path that no implicated file exhibits alone.
The skill collections that supply this ecosystem carry defects at higher rates
than the configurations that consume them, which places the audit at
installation time, and almost half of real configurations run more than one
assistant, which makes vendor-specific tooling structurally insufficient. We
release the tool, the corpus with pinned commit identifiers, and the analysis
pipeline.
\end{abstract}

\begin{IEEEkeywords}
AI code assistants, agent skills, software supply chain security, configuration
auditing, static analysis, LLMOps
\end{IEEEkeywords}

% ===========================================================================
\section{Introduction}
% ===========================================================================

Large language models have become integral to professional software
development. Tools such as Claude Code, GitHub Copilot, and Cursor function as
AI code agents: autonomous systems that combine an LLM with tools, execution
environments, and persistent instructions to perform development tasks on a
developer's behalf. As these agents have been adopted more widely, developers
increasingly customize their behavior through persistent configuration
artifacts, including instruction files, automation hooks, tool integrations, and
subagent definitions that collectively shape how the agent operates within a
project. We refer to this full set of artifacts as the \emph{harness} of the
code agent: the layer that wraps the underlying model and mediates between its
raw capabilities and the developer's intent. The model may be interchangeable
across installations. The harness is what differentiates one developer's agent
from another's.

The harness is a supply chain surface. Its components are routinely installed
from public marketplaces and community repositories, so each installed skill is
an unvetted third-party dependency whose payload is natural-language instruction
and embedded scripts rather than object code. This inverts the usual security
assumption of package management. A conventional dependency is reviewed as code;
a skill is a document that instructs a privileged agent, and it can direct that
agent to read credentials, execute shell commands, or contact the network.
Conventional supply chain defenses, including lockfiles, signature verification,
and CVE scanning of packages, do not inspect this layer. The failure modes are
quiet by construction. A conflicting trigger between two skills throws no error;
it causes unpredictable skill selection. A broken reference to a nonexistent
script fails no build; it silently disables intended behavior. A
prompt-injection pattern embedded in a third-party skill raises no alert; it
creates an attack surface that persists across every session. In each case the
artifact remains syntactically valid and the agent continues to run, which is
precisely why the defect survives.

The risks are empirical rather than hypothetical. A large-scale analysis of the
agent-skill ecosystem \cite{liu2026wild} collected 42{,}447 skills from two
major marketplaces and reported that 26.1\% of the 31{,}132 analyzed carried at
least one vulnerability pattern, with data exfiltration and privilege escalation
most prevalent. In February 2026 a coordinated campaign planted malicious skills
on a public marketplace \cite{yomtov2026clawhavoc}, distributing
credential-stealing malware through install-time instructions. Li et al.
\cite{li2026securing} attribute the most severe threats to structural properties
of the skills framework itself: the absence of a data-instruction boundary, a
single-approval persistent trust model, and the lack of mandatory marketplace
review. CVE-2025-59536 \cite{cve202559536} documents a related mechanism in
which project-scoped configuration takes effect when a repository is opened,
before the developer has consciously approved anything. Drift compounds all of
this: harnesses grow organically, accumulating redundant instructions, broken
references, and components that restate model defaults, and irrelevant context
degrades model reasoning even when the relevant information remains present
\cite{shi2023distracted}, so every wasted token is taken from a useful one.

The difficulty is that existing analyses inspect these artifacts one file at a
time. Per-file inspection answers whether a component is well constructed and
safe, but not whether the assembled harness functions safely as a whole. The
second question requires examining relationships between components: redundancy
across types, contradictory directives, competing demands on the model's
attention, and, critically for security, attacks that become visible only when a
benign-looking component is read against the rest of the setup. Consider two
skills that each individually pass review, one instructing the model to ``always
use \texttt{raise from} for exception chaining'' and another repeating the same
directive. Evaluated independently both are specific and correct, yet loaded
together they double the token cost for one behavioral effect. The same
isolation blindness that misses this redundancy also misses an attack chain in
which a benign command invokes a malicious skill only when a context file
instructs it to. Neither the command nor the skill nor the context file is
anomalous on its own. The chain is the artifact, and we show below that such
chains occur in configurations published today.

Our proposed solution is to make the boundary of per-file analysis explicit and
then audit across it. We organize harness defects by the \emph{scope of
analysis} their detection requires, from a single file up to the full component
graph, and observe that this classification is decidable from a detector's
implementation rather than from its description. This converts a rhetorical
claim, that per-file tools miss things, into a quantity that can be measured
against a real corpus. This paper contributes that organization
(Section~\ref{sec:taxonomy}); an open-source auditor that operationalizes it
through deterministic rules over a unified component graph with reachability
gating, across eight assistant platforms (Section~\ref{sec:tool}); and an
empirical study of 115 assembled configurations and 65 published skill
collections characterizing the defects in each and establishing how much of the
defect surface lies beyond single-file reach (Sections~\ref{sec:design} and
\ref{sec:results}).

% ===========================================================================
\section{Related Work}
% ===========================================================================

\textbf{Agent-skill security.} Liu et al. \cite{liu2026wild} conducted the first
large-scale empirical analysis of the skill ecosystem, combining static analysis
with model-based classification at per-file granularity. Holzbauer et al.
\cite{holzbauer2026context} add repository context to skill classification,
moving toward composition while keeping the individual skill as the unit of
judgment. Schmotz et al. \cite{schmotz2026skillinject} measure agent
susceptibility to skill-file attacks empirically, and Li et al.
\cite{li2026securing} formalize a threat taxonomy validated against confirmed
incidents. Shen et al. \cite{sigil2026} address the gap between what is audited
and what is later loaded at runtime, providing verifiable hosting and load-time
attestation. Li et al. \cite{lesdissonances2025} give the first systematic
security analysis of task control flows in multi-tool agents, identifying
cross-tool harvesting and polluting, and are the closest prior work in spirit,
taking seriously that capability is assembled rather than declared. Closest in
mechanism is the recent proposal of a unified graph representation for agent
security auditing \cite{agentbom2026}, which likewise argues that structural
representation is a precondition for auditing; that work models runtime agentic
systems, whereas we analyze the static configuration artifacts from which such
systems are assembled. These works target marketplace artifacts, runtime
behavior, or distribution integrity. We audit the installed configuration
spanning heterogeneous component types across assistant platforms, and we
compare the consuming population against the supplying one.

\textbf{LLM extension ecosystems.} Analyses of LLM app stores
\cite{hou2025appstores}, ChatGPT plugins \cite{iqbal2024plugins}, and the
ChatGPT app ecosystem \cite{yan2024ecosystem} establish a recurring pattern in
which extension ecosystems accumulate risk faster than review capacity, with the
platform trust model rather than any individual extension usually the root
cause. Indirect prompt injection \cite{greshake2023injection} is the mechanism
by which instruction-as-payload becomes execution. Wang et al.
\cite{wang2025supplychain} frame the LLM supply chain as a research agenda
spanning model, data, and tooling layers; the agent harness is a late addition
to that chain and inherits its characteristic weakness, that installation is
easy and revocation is not.

\textbf{Configuration as an analyzable artifact.} That declarative configuration
warrants systematic analysis has precedent in static analysis of Infrastructure
as Code \cite{chiari2022iac}, catalogs of configuration smells
\cite{sharma2016smell}, and security smells in Puppet \cite{rahman2019sins},
which established that configuration files benefit from structural validation,
security scanning, and cross-file consistency checking. Configuration smells
specific to agent context files \cite{gloaguen2026smells} and an empirical study
of 2{,}303 such files \cite{chatlatanagulchai2025readmes} bring this tradition to
agent configuration, finding among other things that developers rarely specify
security requirements. Both operate at single-file granularity. Related work on
instruction quality establishes why the analysis matters behaviorally: system
prompts materially affect model output \cite{khojah2025prompt}, and prompt
fidelity degrades over multi-turn conversation \cite{li2024instability}. A survey
of agent-harness infrastructure identifies harness-level evaluation as an open
problem \cite{ning2026harness}, and work on validating LLM-assisted evaluation
\cite{shankar2024validators} constrains how far automated verdicts can be
trusted, a caution we adopt in Section~\ref{sec:design}. None of this work
evaluates the composed multi-component configuration as an integrated surface,
which is the gap this paper addresses.

% ===========================================================================
\section{Defects by Analysis Scope}
\label{sec:taxonomy}
% ===========================================================================

Instruction files differ from source code in two ways that make their defects
insidious. They execute probabilistically, influencing model behavior without
guaranteeing it, so a defective instruction produces a distribution of outcomes
rather than a reproducible failure. And they are consumed within a finite
attention budget where every token competes for influence, so a component's cost
is borne by its neighbors. Harness defects therefore produce no error at
configuration time and manifest only at runtime.

Before classifying defects we fix terminology for components, which contemporary
agents expose under different names. \emph{Context files} load into every session
regardless of task (\texttt{CLAUDE.md}, \texttt{AGENTS.md},
\texttt{copilot-instructions.md}), so their cost is unconditional.
\emph{Skills} load when the platform judges them relevant, based on a
description the author writes, so a skill whose description lacks activation cues
may never load and a well-written body can be inert in practice.
\emph{Commands} are user-triggered workflows. \emph{Hooks} are deterministic
scripts bound to lifecycle events, and unlike skills they execute rather than
persuade. \emph{Subagents} are delegated model instances with their own tool
grants. \emph{MCP servers} are external tool providers declared in configuration
and reachable at runtime. A \emph{setup} is a concrete instantiation of some
subset of these for a project, and is distinct from a \emph{collection}, a
repository whose product is a set of publishable skills.

We organize defects by the scope of analysis required to detect them. The useful
property of this axis is that it is decidable from a detector's implementation:
given a rule, one can read what it consumes and determine what scope it needs.

\textbf{FILE.} Decidable from the bytes of one component. Malformed frontmatter,
a hedged instruction that weakens compliance, a prompt-injection pattern, an
over-broad permission grant in a settings file, a hook that leaks environment
variables into its output.

\textbf{FILE\_FS.} Requires the file plus the surrounding filesystem. A
reference to a script that does not exist, taint analysis over a skill's bundled
sources, a CVE lookup against bundled dependencies. The judgment remains local
to one component but cannot be made from its text alone.

\textbf{PAIRWISE.} Requires comparing two components. Content duplication above
a similarity threshold, trigger overlap between skills competing for the same
activation, a command referencing a skill absent from the setup.

\textbf{SETUP.} Requires the whole component graph or an aggregate over all
components. Skills that no component references, circular reference chains,
aggregate context budget against a model's window, permission escalation along
delegation edges, divergence between two assistants' context files, and cross
component flow from a credential-reading component to a network-capable one.

Security is orthogonal and appears at every level, which is why we do not treat
it as a fifth class: a prompt injection is FILE scope, a taint flow through a
bundled script is FILE\_FS, and an attack chain routed through a context file is
SETUP. A per-file scanner is restricted by construction to FILE and to the part
of FILE\_FS its execution model permits. Everything beyond that boundary is what
this paper measures.

The threat model that motivates the SETUP class is worth stating directly. We
consider a developer who installs components from sources they do not fully
control and composes them with components they wrote themselves. The adversary
can publish to a marketplace and rely on installation without code review,
because the artifact reads as documentation; the adversary cannot modify the
agent binary or the model. Under this model, capability is assembled from parts.
A component that reads credentials is unremarkable, a component that reaches the
network is unremarkable, and a delegation edge between them is an exfiltration
path that neither file exhibits alone.

% ===========================================================================
\section{The harness-eval Tool}
\label{sec:tool}
% ===========================================================================

\texttt{harness-eval} is an open-source Python package that discovers harness
components across Claude Code, Cursor, Copilot, Gemini CLI, OpenCode, Windsurf,
Cline, and Codex layouts, parses each into a typed representation carrying
frontmatter, body text, component type, and token count, and reports terminal,
JSON, or SARIF output for code-scanning integration.

\textbf{Deterministic rules.} A pluggable engine ships 97 rules that
self-register through a global registry and are selected by configurable
presets. Each rule emits a structured finding with severity, location, category,
a remediation suggestion, and a mapping to external control frameworks
\cite{owaspllm, mitreatlas}. Rules carry a confidence tier so consumers can gate
a build on precision rather than severity alone, and a baseline mechanism
records existing findings so continuous integration fails only on newly
introduced issues. Deterministic rules are retained rather than delegated to a
model because they execute in milliseconds at zero marginal cost, return
identical results on identical input, operate offline, and cannot fabricate
findings. Determinism is not a stylistic preference: a nondeterministic check
cannot be allowed to break a build intermittently, which rules out model
inference for the gating case.

\textbf{Component graph.} Setup-scope analyses operate over a unified graph
whose nodes are the discovered components and whose edges are the references and
invocations among them. Over this graph the tool computes orphan and circular
reference detection, aggregate context and description budgets against common
model window sizes, permission escalation along delegation edges, divergence
between context files belonging to different assistants, alignment between
configured MCP servers and the skills that reference their tools, and a
cross-component flow analysis pairing a component holding credential or
environment access against a network-capable component reachable from it. We
state the scope of the last analysis precisely: it is a capability-presence and
reachability analysis over graph edges rather than interprocedural taint tracking
across component boundaries. It answers whether a path exists from a
credential-reading component to a network-capable one, not whether a specific
secret reaches a specific sink, which is the appropriate granularity for triage
and the reason it runs in milliseconds. Reachability gating annotates each
security finding with whether the implicated code is reachable from a declared
trigger, so findings in dead code or documentation examples are downgraded and
findings reachable from a broad automatically-loaded trigger are prioritized.

\textbf{Semantic review.} An optional LLM evaluator scores components against
typed rubrics and checks for description-behavior mismatch and semantic attack
variants that evade pattern matching, with every finding required to cite the
content spans it derives from. This layer is deliberately not exercised in the
study below, so that all reported results are reproducible without model access
and free of run-to-run variation.

\textbf{Scope classification of the rule set.} We classified all 97 rules by
reading each rule class and recording what it reads from the analysis context
and the filesystem. An automated pass produced an initial classification, and we
reviewed every rule by hand, correcting 12 whose signals were misleading. Two
patterns recurred: rules reading a settings file appear global because settings
govern the whole project but consume a single component, and rules iterating
over all components while applying an independent per-file check appear
compositional but make each judgment locally. Both are FILE scope. The released
classification records both the automated and corrected results.
Table~\ref{tab:scope} gives the outcome.

\begin{table}[t]
\caption{Rule set by required analysis scope ($n=97$).}
\label{tab:scope}
\centering
\small
\begin{tabular}{lrl}
\toprule
Scope & Rules & Examples \\
\midrule
FILE     & 71 & injection patterns, hedged instructions, \\
         &    & permission grants, hook env leakage \\
FILE\_FS & 10 & broken script references, taint flow, \\
         &    & YARA signatures, CVE lookup \\
PAIRWISE &  6 & duplicate detection, skill overlap, \\
         &    & command referencing an absent skill \\
SETUP    & 10 & orphan skills, circular references, \\
         &    & aggregate budgets, permission escalation, \\
         &    & multi-assistant drift, cross-component flow \\
\midrule
\multicolumn{2}{l}{Beyond single-file scope} & 26 (26.8\%) \\
\bottomrule
\end{tabular}
\end{table}

% ===========================================================================
\section{Study Design}
\label{sec:design}
% ===========================================================================

We ask four questions. \textbf{RQ1}: what do real assembled harness setups look
like, and what defects do they contain? \textbf{RQ2}: how much of the defect
surface lies beyond the reach of single-file analysis? \textbf{RQ3}: how do
assembled setups differ from the published skill collections that supply them?
\textbf{RQ4}: what does auditing cost?

\textbf{Corpus construction.} We built a frame of 5{,}718 unique repositories
using the GitHub repository search API. Fifteen queries covered
assistant-specific topics (\texttt{claude-code}, \texttt{agent-skills},
\texttt{agents-md}, \texttt{cursor-rules}, \texttt{copilot-instructions},
\texttt{opencode}, \texttt{subagents}, and others) together with README
references to canonical context filenames. Each query was issued under two sort
orders, by stars and by recency, over three pages, and deduplicated by
repository name. Sorting by stars alone would bias toward established projects
and miss the long tail of personal configurations where drift accumulates;
sorting by recency alone would bias toward repositories too new to have drifted.
Forks and archived repositories were excluded. From this frame we drew a seeded
random sample and scanned 636 repositories to determine which contain a harness.
Repositories are classified from their component inventory: a \emph{setup} has
at least two distinct component types, or at least one component that is not an
instruction file; a \emph{collection} has five or more skills and nothing that
composes them; the remainder carry a single context file or no components and
are outside the scope of a study about composition. This yields the study
population of 115 setups and 65 collections. The two retained populations are
the two sides of the supply chain, which is why we keep both: collections are
what developers install, setups are what they end up running, and RQ3 asks
whether their defect profiles differ.

\textbf{Scanning protocol.} For each repository the pipeline performed a shallow
single-branch clone, recorded the resolved commit identifier, ran the
deterministic and graph-based analyses under a timeout, recorded findings and
inventory metadata, and deleted the working copy. No third-party repository
content is retained, so the released artifact consists of repository URLs, pinned
commit identifiers, and findings, and redistributes nothing. Because each scan is
pinned to a commit, results are reproducible against the exact tree analyzed even
after the repository changes, which matters in a population this volatile.

\textbf{Rule calibration.} Before the reported run we validated rule behavior
against manual inspection. We re-scanned a stratified sample of repositories,
capturing for every finding the component file, the target line, and a window of
surrounding source, and judged findings against a codebook written before
labeling and released with the artifact. The codebook's governing criterion is
that a finding counts as correct when it identifies a real instance of the
defect the rule names. The validation confirmed that the graph-based rules
carrying this paper's central claims identify real conditions: circular
reference detection was correct in all 11 cases examined, reporting concrete
cycles such as \texttt{prp-new} to \texttt{prp-review} to \texttt{prp-new};
cross-assistant divergence was correct in all 6 cases; unrestricted subagent
grants in all 5. Reference resolution was verified against the filesystem on 113
findings across 10 repositories by re-cloning and testing every flagged path,
and every flagged path failed to resolve, with 76 having no file of that name
anywhere in the repository. Validation also identified three rules whose
heuristics were too coarse for international or deeply structured documents, and
these were revised before the reported run: homoglyph detection now requires
script mixing within a single token, so a document written in Cyrillic is not
flagged for containing Cyrillic while a Latin identifier carrying a substituted
character still is; section-emptiness now treats a subsection as content for its
parent heading; and reference extraction now rejects environment assignments and
normalizes trailing sentence punctuation. All results below were produced with
the calibrated rule set, \texttt{harness-eval} v7.10.1.

\textbf{Scope ablation.} To quantify what per-file tooling reaches, we partition
each repository's findings by the scope of the rule that produced them. The
FILE-only partition is what a per-file scanner operating on identical rule
semantics would report. We prefer this to a head-to-head against a released
scanner, because competing tools implement different rule catalogs, so any
difference in output would confound analysis scope with catalog composition and
would attribute to a competitor a failure to implement rules it never claimed.

% ===========================================================================
\section{Results}
\label{sec:results}
% ===========================================================================

\subsection{RQ1: what real setups contain}

Assembled setups are substantial artifacts. The median setup carries 12
components and 16{,}320 tokens, and the largest carries 207 components and
685{,}356 tokens. Median always-loaded content is 2{,}355 tokens, consumed on
every session regardless of the task at hand.

Multi-assistant configuration is current practice rather than a projection:
48.7\% of setups configure more than one assistant. Claude Code appears in 126
repositories of the study population, OpenCode in 81, Cursor in 16, Copilot in
15, Codex CLI in 13, Gemini CLI in 10, Windsurf in 4, and Cline in 1. A tool
restricted to one vendor's layout would miss part of the configuration in
roughly half the setups it examined.

Defects are common. 85.2\% of setups carry at least one finding and 62.6\% carry
a finding at error severity, meaning a condition that prevents a component from
functioning as written rather than a stylistic concern. Security-category
findings appear in 48.7\% of setups. The most frequent defects are structural
and activation-related. Descriptions lacking activation context appear in 47.8\%
of setups, which is consequential rather than cosmetic: without activation cues
the platform cannot reliably decide when to load a skill, so a carefully written
skill body may never activate for the situation it was written for. Broken
references appear in 44.3\%, silently disabling intended behavior. Unreachable
skills, which no command, context file, or subagent references, appear in 40.9\%.

\subsection{RQ2: the defect surface beyond single-file reach}
\label{sec:rq2}

57.4\% of setups (95\% CI 48.3 to 66.0) carry at least one defect that
single-file analysis cannot reach, and 47.0\% carry one requiring the full
component graph. Measured by finding volume rather than by repository, 32.2\% of
findings in setups come from rules operating beyond file scope
(Table~\ref{tab:ablation}).

\begin{table}[t]
\caption{Findings in assembled setups by required analysis scope.}
\label{tab:ablation}
\centering
\small
\begin{tabular}{lrr}
\toprule
Scope & Findings & Share \\
\midrule
FILE      & 3{,}098 & 67.8\% \\
FILE\_FS  &    725 & 15.9\% \\
SETUP     &    597 & 13.1\% \\
PAIRWISE  &    147 &  3.2\% \\
\midrule
Beyond single file & 1{,}469 & 32.2\% \\
\bottomrule
\end{tabular}
\end{table}

Table~\ref{tab:setupscope} breaks the graph-level findings down by defect.
Unreachable skills dominate by frequency, and this is the rule whose
interpretation depends most on repository type, a point we return to in
Section~\ref{sec:rq3}. Excluding it, 24.3\% of setups still carry a graph-level
defect.

\begin{table}[t]
\caption{Graph-level defects in assembled setups ($n=115$).}
\label{tab:setupscope}
\centering
\small
\begin{tabular}{lrr}
\toprule
Defect & Setups & Share \\
\midrule
Unreachable skills                    & 47 & 40.9\% \\
Cross-assistant divergence            &  9 &  7.8\% \\
Aggregate description budget exceeded &  8 &  7.0\% \\
Circular reference chain              &  7 &  6.1\% \\
Over-permissive delegated grants      &  6 &  5.2\% \\
Aggregate context budget exceeded     &  6 &  5.2\% \\
MCP server with no consuming skill    &  4 &  3.5\% \\
Cross-component credential flow       &  2 &  1.7\% \\
Permission escalation via delegation  &  1 &  0.9\% \\
\bottomrule
\end{tabular}
\end{table}

Three classes deserve comment because, to our knowledge, no published tool
detects them. \emph{Cross-assistant divergence}, in 7.8\% of setups, occurs
where \texttt{CLAUDE.md} and \texttt{AGENTS.md} contain sections that began as
copies and have since diverged, so different agents working on the same project
receive different instructions on the same subject. Given that 48.7\% of setups
run more than one assistant, this class will grow rather than shrink.
\emph{Circular reference chains}, in 6.1\% of setups, are concrete cycles among
named components, for example a command \texttt{prp-new} that routes to
\texttt{prp-review} which routes back to \texttt{prp-new}. Each file is
individually well formed; the defect exists only in the graph.
\emph{Cross-component credential flow}, in 1.7\% of setups, is the class that
motivated this work: a component with credential and environment access holds a
delegation edge to a network-capable component, forming a credential-to-network
path in which neither file, read alone, is remarkable. This is the composed
attack shape of Section~\ref{sec:taxonomy}, observed in configurations published
today rather than constructed for demonstration.

\subsection{RQ3: setups against collections}
\label{sec:rq3}

The two populations differ in ways that bear on where auditing should be applied
(Table~\ref{tab:compare}). Collections carry defects at higher rates on every
dimension, which is the result that matters for the supply chain: the artifacts
developers install are in worse condition than the configurations they assemble.
Two thirds of collections carry a security-category finding, against under half
of setups.

\begin{table}[t]
\caption{Assembled setups against published skill collections.}
\label{tab:compare}
\centering
\small
\begin{tabular}{lrr}
\toprule
 & Setups & Collections \\
 & ($n=115$) & ($n=65$) \\
\midrule
Any finding                     & 85.2\% & 100\% \\
Error-severity finding          & 62.6\% & 80.0\% \\
Security-category finding       & 48.7\% & 66.2\% \\
Beyond single-file scope        & 57.4\% & 100\% \\
Graph-level, excl. reachability & 24.3\% & 35.4\% \\
More than one assistant         & 48.7\% & 26.2\% \\
Median always-loaded tokens     & 2{,}355 & 925 \\
\bottomrule
\end{tabular}
\end{table}

The reachability figure requires care and illustrates why the populations must
be separated. Every collection reports unreachable skills, because a
collection's skills are unreferenced by definition, the skills being the
published product rather than a component of a running configuration. Read
without that distinction, collections would appear to have a universal
composition defect, and a corpus mixing the populations would inflate exactly
the class this paper measures. Excluding reachability, graph-level rates are
35.4\% for collections and 24.3\% for setups, and the comparison becomes
meaningful again. Setups, by contrast, carry more than twice the median
always-loaded token footprint and are twice as likely to run multiple
assistants, so their characteristic risks are budget exhaustion and
cross-assistant inconsistency rather than component-level defects.

\subsection{RQ4: cost}

Median end-to-end cost was 11.9 seconds per repository including clone time, and
the deterministic and graph-based analyses run in pure Python without model
inference, so a full corpus rescan is a matter of minutes rather than dedicated
infrastructure. This matters for deployment: the highest-risk moment for a
harness is not when it is written but when a third-party component is added to
it, and a check that costs seconds can be attached to that event.

% ===========================================================================
\section{Discussion}
% ===========================================================================

\textbf{What the tool audits, and what it does not.} \texttt{harness-eval}
audits whether a harness is structurally sound, internally consistent, and free
of known attack patterns. It does not evaluate whether a configuration improves
developer outcomes. A setup passing every rule may still fail to help if its
instructions are irrelevant to the developer's tasks; a flawed setup may produce
good results when the instructions that load happen to address the developer's
needs. Correctness is not utility and safety is not value: the contract is
triage, surfacing candidates for human attention rather than verdicts on whether
a configuration is good. This echoes the meta-evaluation problem raised by
Shankar et al. \cite{shankar2024validators} and the observation that current
benchmarks cannot cleanly separate the contribution of the harness from that of
the backbone model \cite{ning2026harness}.

\textbf{Structural evidence and textual evidence.} A design lesson emerges from
building and calibrating the tool. Rules whose evidence is a property of the
component graph behave differently from rules whose evidence is a pattern in
prose. A document can contain the words of an attack while describing rather
than performing it, and a security skill legitimately documents the patterns an
attacker would embed, so textual rules in this domain carry an inherent
ambiguity that mitigations reduce without eliminating. A document cannot
accidentally contain a cycle in a reference graph, an unreachable node, or a
delegation edge between a credential reader and a network sink. Where a defect
can be expressed structurally, it should be, and where a textual rule is
unavoidable it benefits from a structural gate such as reachability annotation.
This principle guided calibration and is our main recommendation to others
building detectors for this surface.

\textbf{Threats to validity.} GitHub's code search endpoint, which can match
directly on paths such as \texttt{.claude/skills/}, requires authentication we
did not use, so the frame is topic and README based and over-represents
repositories that advertise their agent tooling; absolute rates describe that
population. Prevalence figures count rule findings: calibration validated the
graph-level rules carrying the central claims and the reference-resolution rule,
but the high-frequency quality rules were not exhaustively validated, and a
single annotator performed the manual portion, so no inter-rater agreement is
reported. Scope was classified from implementation, and 12 of 97 rules required
manual correction, so the axis describes detectors rather than defects in the
abstract, and a different implementation of the same semantic check could fall
in a different class. We report a scope ablation rather than a comparison against
released per-file scanners, which establishes what file-scoped analysis cannot
reach but not that a specific existing tool fails to reach it. Finally, the
corpus is public GitHub repositories; configurations maintained privately inside
organizations may differ systematically and are the population most relevant to
enterprise adoption.

% ===========================================================================
\section{Conclusion}
% ===========================================================================

We treat the AI code-agent configuration as a supply chain surface and a
first-class artifact subject to systematic auditing. Organizing defects by the
scope of analysis their detection requires makes the boundary of per-file
tooling measurable rather than assertable, and measuring it across 115 assembled
configurations shows that the majority carry defects beyond that boundary:
unreachable skills, circular reference chains, divergence between the context
files of different assistants, and credential-to-network delegation paths that
no single file exhibits.

Two results bear on how this surface should be defended. Published skill
collections, the artifacts developers install, carry defects at higher rates than
the configurations developers assemble, which places the audit at installation
time. And almost half of real setups configure more than one assistant, which
makes cross-assistant consistency a present concern and vendor-specific tooling
structurally insufficient. Auditing costs seconds per repository, so it can be
attached to the moment a third-party component enters a harness. We release the
tool, the corpus, and the pipeline to make that auditing reproducible.

% ===========================================================================
\section*{Data Availability}
% ===========================================================================

The tool is available at \url{https://github.com/redhat-community-ai-tools/harness-eval}
and the study artifact, comprising the rule catalog with its scope
classification, the corpus manifest listing every scanned repository with its
pinned commit identifier, all findings, the calibration codebook, and the
scripts that regenerate every number in this paper, at
\url{https://github.com/Benkapner/harness-eval-experiments}. No third-party
repository content is redistributed.

% ===========================================================================
\section*{Ethical Considerations}
% ===========================================================================

All analyzed repositories are public and were evaluated in read-only mode, and
no exploit against any third-party system was developed or run. No previously
undisclosed vulnerability in third-party software was identified: every finding
concerns configuration that is already publicly visible in the repository that
contains it. Repositories are referred to by class rather than by name in the
body of this paper, and identifiers appear only in the released artifact, where
reproducibility requires them. We notify the maintainers of the two setups
exhibiting cross-component credential flow ahead of publication. The study
involved no human subjects and no personally identifiable information.

% ===========================================================================
\begin{thebibliography}{99}
% Ordered by first citation in the text, so numbering follows appearance.
% ===========================================================================

\bibitem{liu2026wild}
Y. Liu, W. Wang, R. Feng, et al., ``Agent Skills in the Wild: An Empirical Study
of Security Vulnerabilities at Scale,'' arXiv:2601.10338, 2026.

\bibitem{yomtov2026clawhavoc}
O. Yomtov, ``ClawHavoc: Malicious Skills on a Public Marketplace,'' Koi Security
research blog, 2026.

\bibitem{li2026securing}
Z. Li, J. Wu, X. Ling, X. Cui, and T. Luo, ``Towards Secure Agent Skills:
Architecture, Threat Taxonomy, and Security Analysis,'' arXiv:2604.02837, 2026.

\bibitem{cve202559536}
``CVE-2025-59536,'' National Vulnerability Database, 2025.

\bibitem{shi2023distracted}
F. Shi, X. Chen, K. Misra, N. Scales, D. Dohan, E. Chi, N. Sch\"arli, and
D. Zhou, ``Large Language Models Can Be Easily Distracted by Irrelevant
Context,'' in \emph{Proc. ICML}, 2023, pp. 31210--31227.

\bibitem{holzbauer2026context}
T. Holzbauer et al., ``Malicious Or Not: Adding Repository Context to Agent
Skill Classification,'' arXiv:2603.16572, 2026.

\bibitem{schmotz2026skillinject}
D. Schmotz, L. Beurer-Kellner, S. Abdelnabi, and M. Andriushchenko,
``Skill-Inject: Measuring Agent Vulnerability to Skill File Attacks,''
arXiv:2602.20156, 2026.

\bibitem{sigil2026}
T. Shen, Y. Feng, K. Zhu, X. Jia, Y. Liu, and L. Zhang, ``Sealing the
Audit-Runtime Gap for LLM Skills,'' arXiv:2605.05274, 2026.

\bibitem{lesdissonances2025}
Z. Li, J. Cui, X. Liao, and L. Xing, ``Les Dissonances: Cross-Tool Harvesting
and Polluting in Pool-of-Tools Empowered LLM Agents,'' in \emph{Proc. Network
and Distributed System Security Symposium (NDSS)}, 2026.

\bibitem{agentbom2026}
``Towards Security-Auditable LLM Agents: A Unified Graph Representation,''
arXiv:2605.06812, 2026.

\bibitem{hou2025appstores}
X. Hou, Y. Zhao, and H. Wang, ``On the (In)Security of LLM App Stores,'' in
\emph{Proc. IEEE Symposium on Security and Privacy}, 2025.

\bibitem{iqbal2024plugins}
U. Iqbal, T. Kohno, and F. Roesner, ``LLM Platform Security: Applying a
Systematic Evaluation Framework to OpenAI's ChatGPT Plugins,'' in \emph{Proc.
AAAI/ACM Conference on AI, Ethics, and Society (AIES)}, 2024.

\bibitem{yan2024ecosystem}
C. Yan, R. Ren, M. Zhang, et al., ``Exploring ChatGPT App Ecosystem:
Distribution, Deployment and Security,'' in \emph{Proc. IEEE/ACM ASE}, 2024.

\bibitem{greshake2023injection}
K. Greshake, S. Abdelnabi, S. Mishra, C. Endres, T. Holz, and M. Fritz, ``Not
What You've Signed Up For: Compromising Real-World LLM-Integrated Applications
with Indirect Prompt Injection,'' in \emph{Proc. ACM Workshop on Artificial
Intelligence and Security (AISec)}, 2023.

\bibitem{wang2025supplychain}
S. Wang, Y. Zhao, X. Hou, and H. Wang, ``Large Language Model Supply Chain: A
Research Agenda,'' \emph{ACM Trans. Softw. Eng. Methodol.}, 2025.

\bibitem{chiari2022iac}
M. Chiari, M. De Pascalis, and M. Pradella, ``Static Analysis of Infrastructure
as Code: A Survey,'' in \emph{Proc. IEEE ICSA-C}, 2022, pp. 218--225.

\bibitem{sharma2016smell}
T. Sharma, M. Fragkoulis, and D. Spinellis, ``Does Your Configuration Code
Smell?'' in \emph{Proc. MSR}, 2016.

\bibitem{rahman2019sins}
A. Rahman, C. Parnin, and L. Williams, ``The Seven Sins: Security Smells in
Infrastructure as Code Scripts,'' in \emph{Proc. ICSE}, 2019.

\bibitem{gloaguen2026smells}
T. Gloaguen et al., ``Configuration Smells in AGENTS.md Files: Common Mistakes
in Configuring Coding Agents,'' arXiv:2606.15828, 2026.

\bibitem{chatlatanagulchai2025readmes}
W. Chatlatanagulchai, H. Li, Y. Kashiwa, et al., ``Agent READMEs: An Empirical
Study of Context Files for Agentic Coding,'' arXiv:2511.12884, 2025.

\bibitem{khojah2025prompt}
R. Khojah, F. Gomes de Oliveira Neto, M. Mohamad, and P. Leitner, ``The Impact
of Prompt Programming on Function-Level Code Generation,'' \emph{IEEE Trans.
Softw. Eng.}, vol. 51, no. 8, 2025.

\bibitem{li2024instability}
K. Li, T. Liu, N. Bashkansky, D. Bau, F. Vi\'egas, H. Pfister, and
M. Wattenberg, ``Measuring and Controlling Instruction (In)Stability in Language
Model Dialogs,'' in \emph{Proc. COLM}, 2024.

\bibitem{ning2026harness}
X. Ning, K. Tieu, D. Fu, et al., ``Code as Agent Harness,'' arXiv:2605.18747,
2026.

\bibitem{shankar2024validators}
S. Shankar, J. D. Zamfirescu-Pereira, B. Hartmann, A. Parameswaran, and
I. Arawjo, ``Who Validates the Validators? Aligning LLM-Assisted Evaluation of
LLM Outputs with Human Preferences,'' in \emph{Proc. ACM UIST}, 2024.

\bibitem{owaspllm}
OWASP, ``OWASP Top 10 for Large Language Model Applications,'' 2025.

\bibitem{mitreatlas}
MITRE, ``MITRE ATLAS: Adversarial Threat Landscape for Artificial-Intelligence
Systems,'' 2025.

\end{thebibliography}

\end{document}
